\begin{figure}[t]
\centering
\begin{tikzpicture}
\begin{axis}[
  width=0.85\linewidth,
  height=5.5cm,
  ybar,
  bar width=14mm,
  ylabel={Held-out KL at 4.5\,bpp},
  symbolic x coords={$L_2$ converged, $L_3$ polish (rolled back), coord descent},
  xtick=data,
  ymin=0.0, ymax=0.55,
  ymajorgrids,
  major grid style={draw=rule!20},
  axis line style={draw=rule},
  tick label style={font=\scriptsize},
  label style={font=\scriptsize},
  title style={font=\scriptsize\itshape},
  title={Qwen3.6-4B 4.5\,bpp anchor: KL across pipeline stages},
  nodes near coords,
  nodes near coords style={font=\scriptsize},
  enlarge x limits=0.25
]
\addplot[draw=accent, fill=accent!40] coordinates {
  ($L_2$ converged, 0.371)
  ($L_3$ polish (rolled back), 0.461)
  (coord descent, 0.245)
};
\end{axis}
\end{tikzpicture}
\caption{Pipeline progression at the Qwen3.6-4B 4.5\,bpp anchor.
$L_2$ converges to KL $0.371$. The $L_3$ polish DP applies the additive-DP
non-additivity bug at scale and \emph{regresses} to KL $0.461$ ($+24\%$);
the real-KL gate detects the regression and rolls the polish output
back to the prior assignment. Coord descent then commits 6 single-Linear
flips of 101 attempted, each gated on real held-out KL, ending at KL
$0.245$ ($-34\%$ versus $L_2$). The pattern is reproducible across every
4B anchor we have measured and is the empirical motivation for separating
candidate generation from candidate selection.}
\label{fig:4b_progression}
\end{figure}
